\section{故障诊断实验与结果分析}
\label{sec:experiments}

\subsection{引言}

前述章节分别完成了汽轮机系统设备健康状态建模、故障知识图谱构建以及KG-GCN故障诊断模型设计。本章在此基础上开展故障诊断算例验证。由于当前机组缺少能够覆盖多类故障的完整现场故障样本，本文采用“真实健康残差背景+独立故障机理约束”的Level-C半物理故障场景进行方法验证。该实验用于检验状态监测残差、知识拓扑和图卷积网络对典型汽轮机故障模式的识别能力，不将半物理构造样本表述为现场真实故障记录。

本文的总体研究对象仍为燃煤机组汽轮机系统，第三章知识图谱覆盖主通流、再热、抽汽回热、给水和冷端等系统。本章选取超高压缸（VHP）作为代表性设备开展工程算例。VHP既包含本体通流状态，又与一级抽汽支路以及一次再热接口直接相连，可在单一设备算例中形成多个具有明确工程位置的故障响应区域，适合用于验证知识图谱约束下的多参数故障诊断方法。

\subsection{研究对象及故障数据集}
\label{subsec:fault_dataset_v3}

\subsubsection{VHP诊断节点}

第二章VHP健康状态模型以16个边界条件和操作变量为输入，对27个候选状态参数进行同步预测。完成状态点筛选后，最终保留24个状态参数作为故障诊断节点。24个节点主要包括VHP叶片级压力、VHP排汽压力与温度、一级抽汽口及逆止门前后参数、1号高压加热器入口抽汽参数、VHP至一次低温再热接口温度以及缸壁和内部蒸汽温度。每一个诊断节点均保持固定DCS身份，并通过第三章设备--参数映射关系投影到VHP局部知识子图。

\subsubsection{故障类型}

结合现有文献故障机理、第三章知识图谱证据以及24个状态参数的可观测范围，构造正常状态和三类典型半物理故障场景，如表~\ref{tab:fault_classes_v3}所示。

\begin{table}[htbp]
    \centering
    \caption{VHP半物理故障诊断类别}
    \label{tab:fault_classes_v3}
    \begin{tabular}{p{1.0cm}p{4.0cm}p{5.1cm}p{4.0cm}}
        \toprule
        类别 & 故障场景 & 主要构造机理 & 主要作用区域 \\
        \midrule
        F0 & 正常状态 & 真实健康残差背景，不叠加故障响应 & 24个状态节点 \\
        F1 & VHP通流部件侵蚀/SPE & 有效通流面积及级组压力分布改变 & VHP本体、排汽及冷再接口 \\
        F2 & 一级抽汽逆止门后管路泄漏 & 抽汽蒸汽泄漏及下游抽汽状态衰减 & 一级抽汽逆止门后及1号高加入口 \\
        F3 & VHP汽封泄漏 & 蒸汽泄漏及本体热力状态改变 & VHP本体、排汽、抽汽及RH1接口 \\
        \bottomrule
    \end{tabular}
\end{table}

需要说明的是，第三章中能够由本机设备资料、原始论文或标准直接支持的关系用于知识图谱构建，而具体24个测点的故障响应方向和幅值仅作为半物理样本构造假设，不反向用于定义KG邻接矩阵，从而避免利用同一故障响应知识同时生成样本和构造分类图所造成的循环论证。

\subsubsection{半物理故障场景构造}

首先利用2024年6月独立测试月DCS实测值与GRU健康模型预测值计算24维健康残差。对第$k$类故障，第$n$个故障残差窗口表示为
\begin{equation}
R_{k,n}^{F}=R_n^{H}+a_{k,n}S_k\odot G_{k,n}+\epsilon_{k,n},
\end{equation}
式中，$R_n^{H}$为真实健康残差窗口，$S_k$为故障空间响应向量，$a_{k,n}$为故障严重程度系数，$G_{k,n}$为故障时间演化函数，$\epsilon_{k,n}$为小幅随机扰动。故障时间过程随机采用阶跃、斜坡或Sigmoid形式，并改变故障起始时刻和局部响应比例，以降低模型对单一固定模板的依赖。

主实验采用60~min残差窗口，故障严重程度随机分布在训练健康残差标准差的$2.0\sim3.0\sigma$范围内。该范围仅用于描述“清晰可观测的典型半物理故障场景”，不对应设备损伤百分比，也不等价于现场故障的绝对严重程度。较弱扰动范围另作为鲁棒性补充实验保留。

\subsubsection{数据划分}

为避免相邻滑动窗口带来的数据泄漏，先按照自然日划分真实健康残差背景，再在各日期内部构造故障窗口。每连续5个自然日中，第1--3日用于训练，第4日用于验证，第5日用于测试；窗口不允许跨日期或跨数据子集。所有标准化参数均仅使用训练日期块计算。

按照上述规则，训练集、验证集和测试集分别包含1712、556和560个窗口，四类状态在各子集中保持均衡，如表~\ref{tab:fault_dataset_split_v3}所示。

\begin{table}[htbp]
    \centering
    \caption{半物理故障数据集样本分布}
    \label{tab:fault_dataset_split_v3}
    \begin{tabular}{lrrrr}
        \toprule
        数据集 & F0 & F1 & F2 & F3 \\
        \midrule
        训练集 & 428 & 428 & 428 & 428 \\
        验证集 & 139 & 139 & 139 & 139 \\
        测试集 & 140 & 140 & 140 & 140 \\
        \bottomrule
    \end{tabular}
\end{table}

\subsection{KG-GCN模型及实验设置}
\label{subsec:experiment_setting_v3}

\subsubsection{模型设置}

为与课题组既有KG-GCN故障诊断研究的实验组织方式保持一致\cite{YangHaodong2026CFB}，设置CNN和AE作为主要对比模型。CNN将$24\times60$残差窗口视为“参数--时间”二维矩阵，通过局部二维卷积核提取故障特征；AE首先以无监督方式重构残差窗口，再通过低维潜变量完成故障分类；KG-GCN则以第三章审计后的VHP参数子图为结构先验，在知识邻接约束下逐层聚合相邻参数残差信息。

KG-GCN采用两层图卷积编码器。首先将每个节点的60点残差序列映射到32维节点表示，然后利用归一化邻接矩阵完成两次邻域传播，并通过残差连接保留节点自身信息。为避免固定DCS节点的工程位置在全局平均池化中被削弱，最后一层保留节点身份，将节点表示按固定次序输入全连接分类层。

\subsubsection{掩码预训练}

考虑现场高质量故障样本较少，本文按照第四章给出的两阶段训练策略实施掩码预训练。训练阶段随机遮蔽15\%的残差元素，遮蔽后的残差矩阵与KG邻接矩阵共同输入GCN编码器，解码器仅对被遮蔽位置计算重构均方误差。完成无标签预训练后冻结图卷积编码器参数，再利用带标签故障窗口训练全连接分类层，并通过Softmax输出故障类别。该技术路线与课题组既有研究中“无标签残差预训练--冻结图编码器--少量标签分类”的方法组织保持一致\cite{YangHaodong2026CFB}，但本文全部模型、数据和文字表述均针对汽轮机系统重新构建。

主要实验参数如表~\ref{tab:kg_gcn_setting_v3}所示。

\begin{table}[htbp]
    \centering
    \caption{KG-GCN主要实验参数}
    \label{tab:kg_gcn_setting_v3}
    \begin{tabular}{ll}
        \toprule
        参数 & 取值 \\
        \midrule
        诊断节点数 & 24 \\
        残差窗口长度 & 60 \\
        GCN层数 & 2 \\
        节点隐藏维数 & 32 \\
        Mask比例 & 15\% \\
        图级读出 & Node-aware flatten \\
        分类隐藏层 & 96 \\
        优化器 & AdamW \\
        学习率 & $1\times10^{-3}\sim1.5\times10^{-3}$ \\
        权重衰减 & $1\times10^{-4}$ \\
        主结果随机种子 & 123、101、202 \\
        \bottomrule
    \end{tabular}
\end{table}

\subsubsection{评价指标}

参照杨昊东KG-GCN算例的评价口径\cite{YangHaodong2026CFB}，本文采用Accuracy、Precision、Recall、F1-score和Macro-F1评价诊断性能，并给出混淆矩阵。对于第$k$类故障，Precision、Recall和F1分别定义为
\begin{equation}
P_k=\frac{TP_k}{TP_k+FP_k},
\end{equation}
\begin{equation}
R_k=\frac{TP_k}{TP_k+FN_k},
\end{equation}
\begin{equation}
F1_k=\frac{2P_kR_k}{P_k+R_k}.
\end{equation}
Macro-F1定义为
\begin{equation}
F1_{macro}=\frac{1}{C}\sum_{k=1}^{C}F1_k,
\end{equation}
整体准确率定义为
\begin{equation}
Accuracy=\frac{N_{correct}}{N_{all}}.
\end{equation}

\subsection{故障诊断结果与分析}
\label{subsec:diagnosis_results_v3}

\subsubsection{不同诊断模型对比}

采用相同的训练、验证和测试日期划分，对CNN、AE和KG-GCN分别使用随机种子123、101和202进行重复训练。结果见表~\ref{tab:main_result_v3}。

\begin{table}[htbp]
    \centering
    \caption{不同故障诊断模型的重复实验结果}
    \label{tab:main_result_v3}
    \begin{tabular}{lcccc}
        \toprule
        模型 & Accuracy/\% & Macro-Precision/\% & Macro-Recall/\% & Macro-F1/\% \\
        \midrule
        AE & $96.43\pm0.82$ & 96.47 & 96.43 & $96.43\pm0.82$ \\
        CNN & $98.33\pm0.45$ & 98.35 & 98.33 & $98.33\pm0.45$ \\
        \textbf{KG-GCN} & $\mathbf{98.45\pm0.10}$ & \textbf{98.48} & \textbf{98.45} & $\mathbf{98.46\pm0.11}$ \\
        \bottomrule
    \end{tabular}
\end{table}

可以看出，三种模型均能够识别清晰可观测的半物理故障场景，其中KG-GCN获得最高平均Accuracy和Macro-F1。虽然在标签充足条件下CNN已接近性能饱和，KG-GCN相对CNN的均值提升较小，但KG-GCN的随机种子标准差明显更低，说明知识拓扑约束下的图特征表示具有更稳定的训练结果。

以seed=123为例，KG-GCN、CNN和AE的误诊样本分别为8、12和25个。KG-GCN混淆矩阵为
\begin{equation}
\begin{bmatrix}
139&1&0&0\\
4&136&0&0\\
1&1&138&0\\
1&0&0&139
\end{bmatrix},
\end{equation}
其中F2一级抽汽支路泄漏和F3 VHP汽封泄漏具有较高的识别精度。F1通流侵蚀与正常状态之间仍存在少量混淆，说明本体通流故障在部分窗口中的残差响应仍可能与真实健康背景重叠。

该现象与知识图谱故障诊断的一般规律一致：当不同故障共享部分局部监测参数时，单纯依赖局部特征容易出现类别交叉，而将设备拓扑和参数关联显式编码到图结构中，可通过邻域信息聚合增强故障传播区域的整体表示\cite{YangHaodong2026CFB}。

\subsubsection{少故障标签条件下的掩码预训练}

为了进一步考察掩码预训练在故障标签稀缺场景中的作用，将带标签训练窗口缩减至20\%。每个随机种子仅保留344个带标签窗口，而无标签掩码预训练仍可以使用全部1712个训练窗口，但完全忽略其故障类别标签。结果见表~\ref{tab:pretrain20_v3}。

\begin{table}[htbp]
    \centering
    \caption{20\%故障标签条件下掩码预训练结果}
    \label{tab:pretrain20_v3}
    \begin{tabular}{lcc}
        \toprule
        训练方式 & Accuracy/\% & Macro-F1/\% \\
        \midrule
        KG-GCN直接监督训练 & $94.46\pm0.94$ & $94.48\pm0.93$ \\
        \textbf{掩码预训练+冻结编码器} & $\mathbf{96.67\pm0.27}$ & $\mathbf{96.67\pm0.27}$ \\
        \bottomrule
    \end{tabular}
\end{table}

掩码预训练使Macro-F1平均提高2.19个百分点，同时随机种子标准差从0.93个百分点降低至0.27个百分点。这表明，当故障标签数量有限时，利用大量无标签运行残差预先学习节点之间的结构关系和正常/异常残差分布，能够有效改善后续小样本故障分类性能。

进一步将带标签窗口减少至5\%，即每个随机种子仅使用84个带标签训练窗口。此时KG-GCN与CNN的结果见表~\ref{tab:small05_v3}。

\begin{table}[htbp]
    \centering
    \caption{5\%故障标签条件下CNN与KG-GCN对比}
    \label{tab:small05_v3}
    \begin{tabular}{lcc}
        \toprule
        模型 & Accuracy/\% & Macro-F1/\% \\
        \midrule
        CNN & $86.85\pm2.20$ & $86.85\pm2.12$ \\
        \textbf{KG-GCN} & $\mathbf{88.04\pm0.31}$ & $\mathbf{88.12\pm0.28}$ \\
        \bottomrule
    \end{tabular}
\end{table}

在极少标签条件下，KG-GCN平均Macro-F1较CNN提高1.27个百分点。更为明显的是，CNN在不同随机种子下的Macro-F1标准差达到2.12个百分点，而KG-GCN仅为0.28个百分点，说明知识图谱与无标签预训练在极少监督样本场景下能够提高模型稳定性。进一步降低至2\%标签时不同方法均出现较大的随机波动，因此本文不将该极端场景作为主要结论。

\subsubsection{知识拓扑消融}

为区分“图卷积结构”与“知识拓扑来源”的贡献，在5\%标签条件下保持掩码预训练、两层GCN和分类头完全一致，仅改变邻接矩阵。Identity图仅保留节点自连接；Corr图按照健康训练残差的绝对相关系数构造统计相关边；KG图采用本机VHP设备拓扑、参数归属及直接热力接口关系。三种图结构结果见表~\ref{tab:graph_ablation05_v3}。

\begin{table}[htbp]
    \centering
    \caption{5\%标签条件下不同图拓扑诊断结果}
    \label{tab:graph_ablation05_v3}
    \begin{tabular}{lcc}
        \toprule
        图拓扑 & Accuracy/\% & Macro-F1/\% \\
        \midrule
        Identity & $86.73\pm0.45$ & $86.78\pm0.41$ \\
        Corr & $87.26\pm1.22$ & $87.33\pm1.11$ \\
        \textbf{KG} & $\mathbf{88.04\pm0.31}$ & $\mathbf{88.12\pm0.28}$ \\
        \bottomrule
    \end{tabular}
\end{table}

KG图相对Identity图和Corr图的Macro-F1分别提高1.34和0.79个百分点，并表现出最低的随机种子波动。该结果说明，当标签充足且各模型接近饱和精度时，图拓扑差异容易被高容量分类器弱化；而在小样本场景下，本机设备拓扑和参数关联提供的先验约束能够减少模型对偶然统计相关性的依赖，知识图谱的贡献更容易显现。

\subsubsection{测点缺失鲁棒性}

考虑DCS测点暂时失效、信号缺失或通信异常等工程情况，进一步在测试阶段保持训练模型不变，随机屏蔽部分状态节点，并以标准化后的零值代替缺失测点。seed=123的代表性结果见表~\ref{tab:missing_random_v3}。

\begin{table}[htbp]
    \centering
    \caption{随机测点缺失条件下的Macro-F1}
    \label{tab:missing_random_v3}
    \begin{tabular}{rcc}
        \toprule
        缺失节点数 & CNN/\% & KG-GCN/\% \\
        \midrule
        0 & 97.85 & \textbf{98.75} \\
        1 & 95.90 & \textbf{97.69} \\
        2 & 94.33 & \textbf{95.74} \\
        4 & 90.29 & \textbf{90.97} \\
        6 & 83.84 & \textbf{84.98} \\
        \bottomrule
    \end{tabular}
\end{table}

随着缺失测点数量增加，两种模型性能均逐步下降，但KG-GCN在各随机缺失比例下均保持更高Macro-F1，表明邻域信息聚合能够在一定程度上减弱单个测点信息丢失的影响。

进一步针对一级抽汽支路构造局部连续测点缺失场景。当仅缺失逆止门后两个温度测点时，CNN的Macro-F1高于KG-GCN；但当逆止门后温度和1号高加入口压力、温度共4个节点同时缺失时，CNN与KG-GCN的Macro-F1分别下降至62.33\%和65.18\%；当逆止门前后及1号高加入口共6个节点同时缺失时，二者分别为58.72\%和62.25\%。此时KG-GCN分别领先约2.85和3.53个百分点，说明在局部物理支路发生较严重信息丢失时，知识拓扑约束下的跨节点传播具有更明显的相对优势。

\subsection{结果讨论}

综合上述实验，本文KG-GCN的优势并不只体现在常规标签条件下略高的平均诊断精度，而表现为一组相互支撑的结果。首先，在标签充足的主实验中，KG-GCN获得最高平均Accuracy和Macro-F1，并具有最低的随机种子波动。其次，在20\%标签条件下，掩码预训练相对直接监督训练获得2.19个百分点的Macro-F1提升，验证了无标签运行残差预训练的有效性。再次，在5\%标签条件下，KG-GCN相对CNN、Identity图和统计相关图均表现出更高平均Macro-F1和更低方差，知识先验的作用更加明显。最后，在随机和局部成组测点缺失条件下，KG-GCN在较严重缺失情况下表现出更好的鲁棒性。

因此，本章结果与第四章提出的方法动机保持一致：利用状态监测模型形成健康残差，在知识图谱约束下组织参数节点关系，再通过掩码预训练和图卷积传播提取多参数耦合故障特征，可以在故障标签稀缺和测点信息不完整的条件下提高诊断稳定性和鲁棒性。

需要再次强调，本章故障样本为基于真实健康残差背景和独立机理知识构造的半物理故障场景，不等同于现场真实故障事件。本文实验结论用于证明方法在机理一致构造场景下的可行性和相对性能，后续仍需结合现场缺陷记录、经校准的热力仿真或真实故障数据进一步验证工程泛化能力。

\subsection{本章小结}

本章以VHP作为燃煤机组汽轮机系统的代表性设备，利用2024年6月真实健康残差背景构造正常状态、通流侵蚀、一级抽汽支路泄漏和VHP汽封泄漏四类半物理故障场景，并采用自然日互斥方式完成训练、验证和测试划分。主实验中KG-GCN取得$98.45\%\pm0.10\%$的Accuracy和$98.46\%\pm0.11\%$的Macro-F1，整体高于CNN和AE。20\%标签条件下，掩码预训练使Macro-F1相对直接监督训练提高2.19个百分点；5\%标签条件下，KG图相对Identity和Corr图分别提高1.34和0.79个百分点。局部抽汽支路4个和6个测点同时缺失时，KG-GCN相对CNN的Macro-F1分别提高约2.85和3.53个百分点。上述结果从诊断精度、小样本稳定性、知识拓扑贡献和测点缺失鲁棒性多个方面验证了所提方法的有效性。
